\subsection{Datasets and tasks}
The experiments use the existing PU, PHM09 and JNU recordings. Table~\ref{tab:data}
summarizes the current protocols. The same tasks and source roles are used
throughout the Q2 comparisons. Each task is repeated with seeds 0--9.

\begin{table}[pos=htbp]
\centering\small
\caption{Current evaluation protocols. Window length and stride are in samples.
Source validation is separated by recording for PU and PHM09 and by temporal
blocks for JNU.}
\label{tab:data}
\begin{tabular}{lrrrrl}
\toprule
Dataset & Classes & Length & Stride & Tasks & Transfer setting\\
\midrule
PU & 3 & 5120 & 4096 & 4 & Three conditions to one\\
PHM09 & 8 & 5333 & 2666 & 20 & Two conditions to one\\
JNU & 4 & 4000 & 2000 & 3 & Two speeds to one\\
\bottomrule
\end{tabular}
\end{table}

PU uses the Paderborn bearing benchmark \cite{pu}. The four recorded conditions
are A: 1500 rpm, 0.1 Nm, 1000 N; B: 1500 rpm, 0.7 Nm, 400 N;
C: 1500 rpm, 0.7 Nm, 1000 N; and D: 900 rpm, 0.7 Nm, 1000 N.
The current inventory contains 18 bearing IDs, each recorded 20 times per
condition at 64 kHz. Labels are the retained healthy, inner-race and outer-race
coarse categories. For each bearing, 16 recordings per source condition are
assigned to training and four to validation using the recorded seed-specific
split. The first 60 complete windows of each recording are used; their last
sample boundary is 246784, so padding is unnecessary. Bearing identities repeat
across conditions. This protocol evaluates transfer to an unseen condition
within the recorded bearing collection. It is different from the historical
six-domain PU protocol that also combines damage-collection membership.

PHM09 uses the gearbox challenge data \cite{phm09}. Its ten domains combine
shaft rates 30, 35, 40, 45 and 50 Hz with High and Low loads: A/B, C/D, E/F, G/H and
I/J, respectively. Sampling frequency is 66666.67 Hz. The eight labels denote
the retained gearbox configurations, with two recording files per
configuration and condition. There are ten adjacent-domain and ten distant-domain
tasks; their exact source--target pairs are listed in Appendix~\ref{app:tasks}.
Recording-grouped source splits retain source-training class coverage.
Configuration labels are diagnostic categories, not gearbox serial identities.

JNU \cite{jnu} has healthy, inner-race, outer-race and rolling-element labels at 600, 800
and 1000 rpm, recorded at 50 kHz. A, B and C denote these speeds. Each source
recording is divided into non-overlapping training and validation time blocks;
window intervals do not cross the split. All three leave-one-speed tasks are
retained. The source recordings are not independent new-bearing replications.
The three leave-one-speed tasks already exhaust the two-source/one-target
combinations of these three speeds. JNU is retained as an additional
cross-dataset applicability case; the non-saturated PU and PHM09 tasks provide
the main discrimination among strong recipes. The existing three-dataset
summary is retained, with per-dataset results reported alongside it.
A separately reported single-source extension evaluates all six cross-speed
directions using source-local numerical analysis (Appendix~\ref{app:jnu-sdg}).

All contexts and paired observations used by the event analyzer lie entirely
inside the corresponding source-training roles. PU pairing retains recorded
bearing identity; PHM09 and JNU use same-class legal source pairs without
inventing an entity ID. Three centered query windows are taken from each
32768-sample source context for the intervention bank. Target data are used
only after all new source models in a comparison are frozen.

\subsection{Training and comparison design}
The matched operator study uses the same convolutional encoder, feature
normalization, classification head and balanced source sampler. The factual
batch size is
$B=DK\lfloor120/(DK)\rfloor$, yielding 117, 112 and 120 for PU, PHM09 and JNU,
respectively. The actual forward uses the $6B$ layout described in
Section~\ref{sec:method}. All matched arms use Adam with learning rate
$10^{-4}$ for 2000 updates and the same initialization and random-draw streams
within each task and seed. Source validation is evaluated at 0, 1000 and 2000;
the final state is primary and earliest strict source best is secondary.
The encoder has three 1D convolutions with 8, 16 and 32 output channels and
kernel lengths 33, 17 and 3. Each convolution is followed by batch
normalization, LeakyReLU with slope 0.1 and average pooling of width 2.
Feature instance normalization, with variance denominator equal to the feature
length and $\epsilon=10^{-5}$, is applied after the final convolutional block.
Flattening gives an 8192-dimensional vector. Fully connected layers of widths
256, 256 and 128 precede the $K$-class output, with LeakyReLU after each hidden
layer. This network is shared by all matched controls.
Table~\ref{tab:algorithm} gives the preparation and learning procedure.
\begin{table}[pos=htbp]
\caption{Source preparation, learning and deployment procedure.}
\label{tab:algorithm}
\begin{minipage}{.98\linewidth}\small
\begin{enumerate}
\item Partition the observed source recordings into training and validation
roles. Select same-class training-context pairs using the declared source
metadata; retain bearing identity for PU.
\item Form the response, shaft-harmonic and stationary dictionaries. Run the
fixed numerical inverse in Eq.~\eqref{eq:inverse} for each registered source pair.
\item For each source context, render $u_1$ and the two fixed-rate $u_q$ signals.
Construct $\mathcal T_q(x)$ with Eq.~\eqref{eq:edit}, then crop the three
observation-supported queries and cache their log-FFT features.
\item For each training update, sample a balanced factual batch and same-class,
same-source-condition bank slots. Draw the shared edit mask and rate direction;
form the auxiliary features using Eq.~\eqref{eq:aux}.
\item Draw the Mixup coefficient and permutation. Update the shared classifier
using Eq.~\eqref{eq:loss} with the common $6B$ forward layout.
\item Freeze the classifier after 2000 updates for the primary evaluation.
Evaluate source-validation selection separately as a secondary policy.
\item At deployment, apply the log-FFT frontend and trained classifier to each
vibration window. Average probabilities only when reporting recording-level
diagnosis.
\end{enumerate}
\end{minipage}
\end{table}

